% Display the course ID while retaining a valid internal problemtools shortname.
\ifdefined\problemheader
  \let\originalproblemheader\problemheader
  \renewcommand{\problemheader}[2]{\originalproblemheader{#1}{problem-36-24527669}}
\fi
\setlength{\parindent}{0pt}
\setlength{\parskip}{4pt}
\problemname{Plate Loading Challenge}

The University Gym is running a plate-loading challenge. Before each
attempt, the instructor present announces a target weight. Your goal is to load the heaviest
barbell that does not exceed the announced target. The equipment officer also wants to
know how many different selections of plate pairs could produce that desired weight.

The empty bar alone weighs $B$ units. The rack holds $n$ labelled pairs of plates,
numbered $1$ to $n$. Pair $i$ contains two plates, each weighing $p_i$ units.
To use a pair, load each end of the bar with the plates, adding $2p_i$
units in total. Each pair may only be used at most once, and it cannot be split up or combined 
with another. There is also enough room on the bar for any number of selections.

2 selections are different if they use different sets of pair labels.
Pairs with equal weights still have different labels. 
Rearranging the selected plates does not create a different selection. 
Selecting no pairs is allowed and it also counts as 1 selection.

All weights are expressed in integer units, where one unit is =  $100$ grams.
The target $T$ includes the weight of the bar. Find the maximum achievable
total weight at most $T$, and the exact number of selections achieving it.

\section*{Input}
The first line contains three integers $n$, $B$, and $T$, where
$1 \leq n \leq 200$, $1 \leq B \leq 1000$, and $B \leq T \leq B+10000$.

The second line contains $n$ integers $p_1,\ldots,p_n$, where
$1 \leq p_i \leq 1000$. The $i$th integer is the weight of \emph{one plate}
from pair $i$, in the same units as $B$ and $T$.
There is one challenge per input file.

\section*{Output}
Output two integers on one line, separated by a space: the maximum total
barbell weight in the given units, and the exact number of selections
achieving that weight.
The count may exceed the range of a 64-bit integer.

\section*{Sample explanations}
In Sample 1, the empty bar weighs $200$ units ($20$ kg).
The selections $\{1,2\}$, $\{1,3\}$, and $\{2,3,4\}$ each add $160$
units. They are the three selections achieving the target of $360$ units.

In Sample 2, pair 1 alone gives $212$ units and pair 2 alone gives $220$
units. Both pairs together would exceed the target of $221$ units, so the
answer is $220$ units with one selection.

In Sample 3, the target equals the empty bar's weight. The only possible
selection is the empty selection.
